\documentclass[../../main.tex]{subfiles}% 注意这里的文件路径不能用 ./main.tex ,否则用latexmk编译子文件会报错
\graphicspath{{\subfix{./image/}}{\subfix{../../image/}}} % 指定图片引用路径,后续可以直接使用图片文件名
% 如果图片存放在上级目录,则使用 ../../image/ 作为备用路径

% 例如:
% \begin{figure}[H]
% \centering
% \includegraphics[width=0.3\textwidth]{图.png}
% \caption{}
% \label{figure:图}
% \end{figure}
% 注意:上述\label{}一定要放在\caption{}之后,否则引用图片序号会只会显示??.

\begin{document}

\section{偏微分方程的基本概念}

\begin{definition}
\textbf{一个偏微分方程}是与一个未知的多元函数及它的偏导数有关的方程；\textbf{一个偏微分方程组}是与多个未知的多元函数及它们的偏导数有关的方程组。
\end{definition}

\begin{definition}
设 $\Omega\subset \mathbb{R}^n$，$x=(x_1,x_2,\cdots,x_n)$ 表示 $\Omega$ 上的点。假设 $u=u(x):\Omega\to\mathbb{R}$ 是一个函数。对固定正整数 $k$，我们用符号 $D^k u$ 表示 $u$ 的所有 $k$ 阶偏导数
\begin{align*}
\frac{\partial^k u}{\partial x_{i_1}\partial x_{i_2}\cdots\partial x_{i_k}},
\end{align*}
其中 $(i_1,i_2,\cdots,i_k)$ 是集合 $\{1,2,\cdots,n\}$ 中 $k$ 个元素的任意排列。因此，我们可以把 $D^k u$ 看成是 $n^k$ 维欧氏空间 $\mathbb{R}^{n^k}$ 上的向量，并记它的长度为
\begin{align*}
|D^k u| = \left( \sum_{i_1=1}^n \sum_{i_2=1}^n \cdots \sum_{i_k=1}^n \left| \frac{\partial^k u}{\partial x_{i_1}\partial x_{i_2}\cdots\partial x_{i_k}} \right|^2 \right)^{\frac{1}{2}}.
\end{align*}

特别地，当 $k=1$ 时，我们称 $n$ 维向量
\begin{align*}
Du  = \nabla u = \left( \frac{\partial u}{\partial x_1}, \frac{\partial u}{\partial x_2}, \cdots, \frac{\partial u}{\partial x_n} \right)
\end{align*}
为 $u$ 的\textbf{梯度}；当 $k=2$ 时，我们称 $n\times n$ 矩阵
\begin{align*}
D^2 u = \begin{pmatrix}
\frac{\partial^2 u}{\partial x_1^2} & \frac{\partial^2 u}{\partial x_1\partial x_2} & \cdots & \frac{\partial^2 u}{\partial x_1\partial x_n} \\
\frac{\partial^2 u}{\partial x_2\partial x_1} & \frac{\partial^2 u}{\partial x_2^2} & \cdots & \frac{\partial^2 u}{\partial x_2\partial x_n} \\
\vdots & \vdots & \ddots & \vdots \\
\frac{\partial^2 u}{\partial x_n\partial x_1} & \frac{\partial^2 u}{\partial x_n\partial x_2} & \cdots & \frac{\partial^2 u}{\partial x_n^2}
\end{pmatrix}
\end{align*}
为 $u$ 的 $\mathbf{Hessian}$ \textbf{矩阵}。通常记符号 $\Delta$ 为 $\mathbf{Laplace}$\textbf{算子}，
\begin{align*}
\Delta u = \operatorname{tr}(D^2 u) = \sum_{i=1}^n \frac{\partial^2 u}{\partial x_i^2},
\end{align*}
也就是 $u$ 的 Hessian 矩阵的迹，即 $u$ 的 Hessian 矩阵的对角线元素之和。
\end{definition}

\begin{definition}
设 $\Omega\subset \mathbb{R}^n$,$F=(F_1,F_2,\cdots,F_n):\Omega\to\mathbb{R}^n$ 是一个向量函数，记 $F$ 的\textbf{散度}为
\begin{align*}
\operatorname{div} F = \sum_{i=1}^n \frac{\partial F_i}{\partial x_i}.
\end{align*}
于是 $\Delta u$ 是 $u$ 的梯度的散度，即
\begin{align*}
\Delta u = \operatorname{div}(Du).
\end{align*}
为简单起见，我们通常也使用如下符号
\begin{align*}
u_{x_i} = \frac{\partial u}{\partial x_i}, \qquad u_{x_i x_j} = \frac{\partial^2 u}{\partial x_i\partial x_j}.
\end{align*}
\end{definition}

\begin{definition}
设 $\Omega\subset \mathbb{R}^n$,我们用 $C(\Omega)$ 表示在 $\Omega$ 上的连续函数构成的线性空间。对于它的元素 $u\in C(\Omega)$，其模定义为
\begin{align*}
\|u\|_{C(\Omega)} = \sup_{x\in\Omega} |u(x)|.
\end{align*}
我们用 $C^k(\Omega)$ 表示 $\Omega$ 上所有 $k$ 阶偏导数都存在和连续的函数构成的线性空间，也就是在 $\Omega$ 上 $k$ 次连续可微的函数构成的线性空间。对于它的元素 $u\in C^k(\Omega)$，其模定义为
\begin{align*}
\|u\|_{C^k(\Omega)} = \sup_{x\in\Omega} |u(x)| + \sum_{|\alpha|=1}^k \sup_{x\in\Omega} |D^\alpha u(x)|,
\end{align*}
其中 $\alpha=(\alpha_1,\alpha_2,\cdots,\alpha_n)$ 表示多重指标，
\begin{align*}
|\alpha| = \alpha_1+\alpha_2+\cdots+\alpha_n, \qquad D^\alpha u = \frac{\partial^{|\alpha|} u}{\partial x_1^{\alpha_1}\partial x_2^{\alpha_2}\cdots\partial x_n^{\alpha_n}}.
\end{align*}
对于函数 $u\in C(\Omega)$，定义 $u$ 的支集为所有满足 $u(x)\neq 0$ 的点集在 $\Omega$ 上的闭包，记为
\begin{align*}
\operatorname{spt} u = \overline{\{x\in\Omega \mid u(x)\neq 0\}}.
\end{align*}
我们用 $C_0^k(\Omega)$ 表示 $C^k(\Omega)$ 中具有紧支集的函数类。同时我们用 $C^\infty(\Omega)$ 表示在 $\Omega$ 上任意阶偏导数都存在和连续的函数类，即
\begin{align*}
C^\infty(\Omega) = \bigcap_{k=1}^\infty C^k(\Omega).
\end{align*}
\end{definition}


\begin{definition}
设 $\Omega\subset \mathbb{R}^n$,如下形式的方程
\begin{align}
F[D^k u(x), D^{k-1}u(x), \cdots, Du(x), u(x), x] = 0,\quad x\in\Omega. \label{eq:pde-general}
\end{align}
称为一个 $\boldsymbol{k}$ \textbf{阶偏微分方程}，其中
\begin{align*}
F:\mathbb{R}^{n^k}\times\mathbb{R}^{n^{k-1}}\times\cdots\times\mathbb{R}^n\times\mathbb{R}\times\Omega\to\mathbb{R}
\end{align*}
是一个给定函数，$u:\Omega\to\mathbb{R}$ 是一个未知函数。一个偏微分方程的阶就是此偏微分方程中出现的未知函数的偏导数的最高次数。

我们将满足方程 \eqref{eq:pde-general} 的所有函数称为方程 \eqref{eq:pde-general} 的解。
如果 $u\in C^k(\Omega)$ 满足方程 \eqref{eq:pde-general}，则我们称它是方程 \eqref{eq:pde-general} 的\textbf{古典解}。
\begin{enumerate}[(1)]
\item 如果方程 \eqref{eq:pde-general} 可表示成
\begin{align*}
\sum_{|\alpha|\leqslant k} a_{\alpha}(x) D^{\alpha} u = f(x),
\end{align*}
其中 $a_{\alpha}$ $(|\alpha|\leqslant k)$ 和 $f$ 是给定的函数，则称方程 \eqref{eq:pde-general} 为\textbf{线性偏微分方程}；
\item 如果方程 \eqref{eq:pde-general} 可表示成
\begin{align*}
\sum_{|\alpha|=k} a_{\alpha}(x) D^{\alpha} u = f[D^{k-1}u(x), \cdots, Du(x), u(x), x],
\end{align*}
其中 $a_{\alpha}$ $(|\alpha|=k)$ 和 $f$ 是给定函数，则称方程 \eqref{eq:pde-general} 为\textbf{半线性偏微分方程}；
\item 如果方程 \eqref{eq:pde-general} 可表示成
\begin{align*}
\sum_{|\alpha|=k} a_{\alpha}[D^{k-1}u(x), \cdots, Du(x), u(x), x] D^{\alpha} u
= f[D^{k-1}u(x), \cdots, Du(x), u(x), x],
\end{align*}
其中 $a_{\alpha}$ ($|\alpha|=k$) 和 $f$ 是给定函数，则称方程 \eqref{eq:pde-general} 为\textbf{拟线性偏微分方程}；
\item 如果方程 \eqref{eq:pde-general} 非线性地依赖于 $u(x)$ 的最高阶偏导数 $D^k u$，则称方程 \eqref{eq:pde-general} 为\textbf{完全非线性偏微分方程}。
\end{enumerate}
\end{definition}







\end{document}